\documentclass[11pt,letterpaper]{article}

% Margins
\usepackage[left=0.75in,right=0.75in,top=0.75in,bottom=0.75in]{geometry}

% Font
\usepackage[T1]{fontenc}
\usepackage{lmodern}
\usepackage{bbm}
\usepackage{dsfont}


% Hyperlinks
\usepackage{hyperref}
\hypersetup{
    colorlinks=true,
    urlcolor=blue,
    linkcolor=blue
}

% No page numbers
\pagestyle{empty}

% Spacing
\usepackage{parskip}
\setlength{\parindent}{0pt}

% Section formatting
\usepackage{titlesec}
\titleformat{\section}{\large\bfseries}{}{0em}{}[\titlerule]
\titlespacing*{\section}{0pt}{1.5ex}{0.5ex}

% Custom commands
\newcommand{\dateright}[1]{\hfill\textnormal{#1}}
\newcommand{\entry}[4]{%
    \textbf{#1} \dateright{#3}\\
    #2 \dateright{#4}
}

\begin{document}

% Header
{\LARGE\textbf{Sabyasachi Mukherjee, M.Sc.}}

\vspace{0.2cm}

Email: \href{mailto:sabyasachi.mukherjee@smail.uni-koeln.de}{sabyasachi.mukherjee(at)smail(dot)uni-koeln.de} \hspace{0.3cm} Website: \href{https://www.mi.uni-koeln.de/~smukherj/}{https://www.mi.uni-koeln.de/~smukherj/}




\section*{Interests}

Analytic number theory, modular forms and packing problems

\section*{Education}

\entry{PhD in Mathematics}{Universität zu Köln}{April 2026 - Present}{Cologne, Germany}

\entry{M.Sc. in Mathematics}{Georg-August-Universität Göttingen}{March 2022}{Göttingen, Germany}\\
Thesis: \textit{The Möbius Function on Arithmetic Progressions}

\vspace{0.2cm}

\entry{B.S. in Mathematics}{Shiv Nadar University}{May 2017}{Noida, India}

%Thesis: \textit{Dirichlet's Theorem on Arithmetic Progressions}

\vspace{0.2cm}


\section*{Teaching \& Grading Experience}

\entry{Tutor}{Real Analysis I}{Georg-August-Universität Göttingen}{Winter 2020}
%(Differenzial- und Integralrechnung I)
\vspace{0.2cm}

\entry{Tutor}{Functional Analysis}{Georg-August-Universität Göttingen}{Summer 2020}

\vspace{0.2cm}

\entry{Tutor}{Real Analysis I }{Georg-August-Universität Göttingen}{Winter 2019}
%(Differenzial- und Integralrechnung I)
\vspace{0.2cm}

\entry{Grader}{Linear Algebra and Analytic Geometry}{Georg-August-Universität Göttingen}{Summer 2019}

\section*{Selected Talks}

\textbf{Hasse Principle for Non-singular Quartic Hypersurfaces} \dateright{February 2022}\\
Analytic Number Theory Working Seminar

\vspace{0.3cm}

\textbf{A Sharp Version of Halász's Theorem} \dateright{December 2020}\\
Göttingen Number Theory Working Seminar


\vspace{0.3cm}


\textbf{Large subsets of $\mathds{F}_q^n$ without three-term arithmetic progressions} \dateright{July 2020}\\
Göttingen Number Theory Oberseminar



\section*{Conferences \& Summer Schools}

\textbf{Summer School in Analytic Number Theory} \dateright{June 28 -- July 2, 2021}\\
Paris, France

\vspace{0.1cm}

\textbf{Hausdorff Centre of Mathematics: "The Circle Method"} \dateright{May 10--21, 2021}\\
Bonn, Germany





\section*{Technical Skills and Languages}

Python, MATLAB, \LaTeX, English (C2), German (C1), Bengali (C2), Hindi (C2)


\section*{Work Experience}

\textbf{Industry}: Data Science \dateright{2022 -- 2025}


\end{document}